\section{Parasitics in the Installed Magnetics}
\label{sec:mag-parasitics}

The resonant path includes the added series inductor, transformer leakage, and connecting conductors. The following first-order relationship was used to check how installation changes resonance:
\begin{equation}
 f_{r,\mathrm{installed}}\simeq
 \frac{1}{2\pi\sqrt{(L_{\mathrm{RES}}+L_{\mathrm{lk}}+L_{\mathrm{int}})C_{\mathrm{RES}}}},
 \label{eq:mag-installed-resonance}
\end{equation}
Here \(L_{\mathrm{int}}\) is the added connection inductance. Both component values and wiring can therefore change tank current and commutation timing. Magnetizing inductance, load, and nonlinear device capacitance must also be included in a more complete model.

\begin{table}[htbp]
\centering
\caption{Parasitic effects in the installed magnetics.}
\label{tab:mag-parasitic-mechanisms}
\begin{tabularx}{\textwidth}{@{}p{0.28\textwidth}X@{}}
\toprule
Mechanism & Consequence for the converter\\
\midrule
Connection inductance & Resonance shift, switching overshoot, and ringing\\
Connection resistance & Additional loss and local heating at pads and joints\\
Primary-to-secondary capacitance & Common-mode current driven by the switch-node transition\\
Core-gap variation & Change in inductance and magnetizing current\\
Polarity or terminal error & Incorrect voltage transfer or rectifier operation\\
\bottomrule
\end{tabularx}
\end{table}

Common-mode current follows \(i_{\mathrm{CM}}=C_{\mathrm{ps}}\,dv/dt\). Interwinding capacitance, adjacent copper, and probe loading therefore affect secondary-noise captures. Quantifying these effects requires impedance and capacitance measurements on the assembled converter.
